\section{The residue quotient and the finite stable image}
\label{sec:cutoff}

This section has two outputs with different hypotheses. First, we prove
that the stable Cartier image on the cohomology of a compact coefficient
lattice is finite. The finite Laurent homology of
\cref{cor:laurent-finite} supplies this result through a residue
quotient and compact duality. No next-height finiteness hypothesis is
needed for this step.

The second output concerns the corresponding intersection in localized
cohomology. Here we assume next-height finiteness and countability of
one consecutive quotient of Cartier images. These additional inputs
make the localization kernel finite on the relevant compact image and
allow us to identify the two intersections. Section~\ref{sec:pro-transfer}
will establish the consecutive-image hypothesis and use the resulting
finite localized intersection.

Fix a height stratum and a sufficiently deep finite marking level as in
\cref{sec:transfers,sec:laurent}.  All vector spaces in this section are
over the \emph{finite} field $\kk$.  Enlarge $\kk$ and the fixed power
$h_0=p^{\nu_0}>1$, if necessary, so that $h_0$-power Frobenius fixes $\kk$.
Write
\[
 A=\kk[[x,y_1,\ldots,y_{r-1}]],\qquad
 B'=A'[1/x],\qquad B_U=eB',
\]
where $A'$ is the whole finite free model, with a fixed integral basis
$b_1,\ldots,b_t$, and $e$ is its $G$-invariant idempotent on the open.
We use the normalized $\kk$-linear Cartier transfer $S$ of
\cref{prop:trans-cartier}; on $B'$ it means $f\mapsto S(ef)$.
Thus its output lies in $eB'$, and it is zero on $(1-e)B'$.

There is a constant $c_S\geq0$, independent of the nonnegative integer
$D$, such that
\begin{equation}
 S(x^{-D}A')\subseteq
 x^{-\lceil D/h_0\rceil-c_S}A'.
 \label{eq:cutoff-pole-recurrence}
\end{equation}
All coefficient maps are $G$-equivariant and are continuous on a
bounded pole lattice into the indicated lattice.  In particular they
act on the cochain and chain complexes of \cref{sec:continuous}.
Choose an integer $D_0$ sufficiently large that
\begin{equation}
 D_0>\frac{c_S+1}{1-h_0^{-1}},\qquad
 \mathcal P=x^{-D_0}A'.
 \label{eq:cutoff-lattice-choice}
\end{equation}
We will increase $D_0$ once more for a residue bound below.  The lattice
$\mathcal P$ is compact, metrizable, $G$-stable, and $S$-stable.
Moreover, for every integer $D\geq0$ there exists $a_0(D)$ such that
\begin{equation}
 S^a(x^{-D}A')\subseteq\mathcal P
 \quad\text{for all }a\geq a_0(D).
 \label{eq:cutoff-eventual-integrality}
\end{equation}
Indeed iteration of \cref{eq:cutoff-pole-recurrence}, with the rounding
absorbed by $c_S+1$, gives the upper bound
\[
 h_0^{-a}D+(c_S+1)\frac{1-h_0^{-a}}{1-h_0^{-1}}
\]
for the pole order.  The strict inequality in
\cref{eq:cutoff-lattice-choice} proves
\cref{eq:cutoff-eventual-integrality}.  This conclusion is uniform on
each indicated lattice and on compact families of its elements; its
exponent is allowed to depend on $D$.

\subsection{Normalized bases and the positive Laurent subspace}

Put $h=h_0^a$, $x_h=x^{1/h}$, $y_{j,h}=y_j^{1/h}$, and
$b_{i,h}=b_i^{1/h}$.  We regard all these root rings inside a common
perfected coefficient ring.  The $b_{i,h}$ are the
\emph{normalized basis at level $h$}; they form a basis of
$B'^{1/h}$ over $B^{1/h}$ and of the corresponding analytic coefficient
module over the perfected base.  This is an actual root basis, not a
new choice of monomials made separately for each element.

The direction of the change of basis is important.  If $h'=h h_0^v$,
then $b_{i,h}$ expands on the $b_{j,h'}$ with coefficients in
$A^{1/h'}$.  To see this, expand $b_i^{h_0^v}$ on the integral basis of
$A'$ and take its $h'$th roots.  Consequently this change of basis
introduces neither negative x-orders nor negative y-orders.  The
opposite change of basis can have x-poles, but at each fixed level its
matrix has bounded x-poles and no negative y-orders.  The topology and
cofinal complete steps supplied by \cref{prop:laurent-model} are
independent of these fixed changes of basis.

Let $\eta$ be the invariant volume of \cref{prop:trans-volume}.  Extend
it arbitrarily by a volume generator on the complementary component of
$B'$.  At level $h$ write its transported version as
\[
 \eta_h=\theta_h\,d\log x_h\wedge
       dy_{1,h}\wedge\cdots\wedge dy_{r-1,h},
 \qquad \theta_h\in (B'^{1/h})^\times.
\]
The extension on the complementary component need not be invariant;
all expressions eventually paired will lie in the $e$-component.
Finite étaleness makes the trace pairing perfect.  The matrix
\begin{equation}
 \mathcal T_h=
 \bigl(\Tr_{B'^{1/h}/B^{1/h}}
             (b_{i,h}b_{j,h}\theta_h)\bigr)_{i,j}
 \label{eq:cutoff-trace-matrix}
\end{equation}
and its inverse act on the negative-y Laurent quotient.
By increasing a fixed integer $c\geq1$, we may assume that their
coefficients, and all fixed multiplication matrices used with them,
have x-poles at most $c/h$ and no negative y-orders.  In fact these are
the root transports of fixed matrices at level one.  Increase $D_0$ so
that
\begin{equation}
 D_0>c.
 \label{eq:cutoff-residue-margin}
\end{equation}
The choice still satisfies \cref{eq:cutoff-lattice-choice}.

For a compact $\kk$-vector space $V$ write
$V^\vee=\Hom_{\kk,\cts}(V,\kk)$ with the discrete topology.
Composing with $\Tr_{\kk/\Fp}$ identifies this with the usual
finite-field Pontryagin dual.  Define
\begin{equation}
 Q=\colim\bigl(
       \mathcal P^\vee\xrightarrow{S^\vee}\mathcal P^\vee
       \xrightarrow{S^\vee}\cdots\bigr),
 \label{eq:cutoff-Q}
\end{equation}
and give $Q$ the discrete topology.  The dual of $\mathcal P$ is
countable: a continuous functional depends on only finitely many
coefficients in its whole finite basis, and $\kk$ and the monomial
index set are finite and countable, respectively.  Thus $Q$ is
countable.  The actions on these discrete duals are continuous, since
the original action on the compact lattice is jointly continuous.

We recall explicitly the Laurent convention from
\cref{prop:laurent-model}.  An element of $W$ is an $e$-supported
expression on the whole basis with monomials
\[
 x^I y_1^{-J_1}\cdots y_{r-1}^{-J_{r-1}},
 \qquad I\in\Z[1/p],\quad
 J_j\in\Z[1/p]_{>0},
\]
whose coefficients are in $\kk$ and tend to zero in some weighted
complete step
\begin{equation}
 w_{L,\delta}(I,J)=
 \begin{cases}
 LI-\delta|J|,&I\geq0,\\
 \delta I-\delta|J|,&I<0,
 \end{cases}
 \quad L>\delta>0,
 \qquad |J|=\sum_jJ_j.
 \label{eq:cutoff-weight}
\end{equation}
The norm of a nonzero monomial is $\exp(-w_{L,\delta}(I,J))$.
When $r=1$ there are no y-variables, and all references to their
negative parts are omitted.  Multiplication by an analytic coefficient
is followed by projection to the part strictly negative in every
y-variable.  Root denominators in an element of $W$ need not be
bounded.

\begin{definition}
An element $w\in W$ belongs to $W_+$ if, in some normalized basis
$b_{i,h}$, every nonzero coefficient monomial of $w$ has
$I\geq\epsilon$ for some real $\epsilon>0$.  The bound is uniform
over its finitely many basis components.  This condition imposes no
bound on the denominators of the exponents in those components.
\end{definition}

Positivity persists on every later normalized basis, by the integral
change of basis described above.  Hence the definition makes $W_+$ a
vector subspace: two elements may be compared at one later level and
with the smaller of their two positive bounds.  The absolute
Frobenius operator
\[
 \Phi:W\longrightarrow W,\qquad w\longmapsto w^{h_0},
\]
is a $\kk$-linear automorphism, as established on the support model
in \cref{prop:laurent-model}.  It scales all exponents and transports
the actual root basis.  Both $\Phi$ and $\Phi^{-1}$ preserve $W_+$.
This operator is distinct from the coefficient-only Frobenius used
to obtain $W$ in \cref{eq:laurent-descent}.

Our aim is to identify $Q$ with $W/W_+$. At root level $h$, the
compact test lattice and the trace matrices have pole bounds of order
$1/h$. Hence a Laurent expression with a fixed positive x-gap is
invisible to all sufficiently late residue tests. Conversely, a series
with finite-field coefficients has only finitely many nonzero terms
below any fixed x-order. Truncating those terms will define the residue
class even when the remaining series has unbounded root denominators.

After constructing this quotient, we will prove that $W_+$ has zero
group homology. Absolute Frobenius supplies the contraction, and
countability will make boundaries open in each complete cycle space.
Together these facts turn the contraction into an actual boundary
identity.

\subsection{The residue map and its kernel}

Let $W_{\mathrm{fin}}$ be the subspace of elements whose root
denominators are bounded; equivalently, the element can be expressed
in a normalized basis at some level $h$ with all exponents in
$h^{-1}\Z$.  It can still be an infinite convergent series.  For
$w\in W_{\mathrm{fin}}$ at this level and
$p_h\in\mathcal P^{1/h}$ define
\begin{equation}
 \lambda_h(w)(p_h)=
 \res_{x_h,y_h}
 \Tr_{B'^{1/h}/B^{1/h}}(w p_h\eta_h).
 \label{eq:cutoff-residue-pairing}
\end{equation}
In logarithmic x-coordinates the residue extracts x-exponent zero and
y-exponent $-1$ in every $y_{j,h}$.  This is also the ordinary formal
residue after replacing $d\log x_h$ by $dx_h/x_h$.

There are only finitely many nonzero terms with $I\leq M$ in any
$\kk$-series satisfying \cref{eq:cutoff-weight}, for every real $M$.
Indeed their weights are bounded above by $\max(LM,0)$, whereas a
weighted null series has only finitely many terms below any given
weight bound.  This uses finiteness of $\kk$: a nonzero coefficient
has norm one.  Applying the observation with the pole bounds for
$p_h$ and $\mathcal T_h$ shows that
\cref{eq:cutoff-residue-pairing} uses only finitely many terms of $w$
and finitely many coefficients of $p_h$.  It is therefore a
well-defined continuous functional on $\mathcal P^{1/h}$.

Powering identifies $\mathcal P^{1/h}$ with $\mathcal P$ as a
compact $\kk$-vector space.  Under this identification the map
between consecutive compact lattices is
\[
 T_h:\mathcal P^{1/(h h_0)}\longrightarrow\mathcal P^{1/h},
 \qquad T_h(v)=\bigl(S(v^{h h_0})\bigr)^{1/h}.
\]
The Cartier-residue identity gives
\begin{equation}
 \lambda_{h h_0}(w)=T_h^\vee\lambda_h(w).
 \label{eq:cutoff-residue-transition}
\end{equation}
For completeness, in one ordinary coordinate with $\eta=dx$, Cartier
takes $x^j$ to $x^{(j+1)/h_0-1}$ when the displayed exponent is an
integer and to zero otherwise.  Pairing this with $x^I$ gives the same
coefficient as pairing $x^I$ at the next normalized root level with
the original monomial.  With $d\log x$ the condition is divisibility
of $j$ by $h_0$ instead.  Taking products of these coefficient
identities proves the assertion in several coordinates.  Trace
commutes with the étale root transport, and multiplication by
$\theta_h$ gives \cref{eq:cutoff-residue-transition} for the volume
used here; this is the compatibility in
\cref{prop:trans-cartier}.  Thus the $\lambda_h$ define a map
$\pi_{\mathrm{fin}}:W_{\mathrm{fin}}\to Q$.

\begin{proposition}[The discrete residue quotient]
\label{prop:cutoff-residue}
The finite-level pairings extend uniquely to an exact sequence of
$\kk[G]$-modules
\begin{equation}
 0\longrightarrow W_+\longrightarrow W
   \xrightarrow{\ \pi\ }Q\longrightarrow0.
 \label{eq:cutoff-residue-sequence}
\end{equation}
The map $\pi$ is continuous on every complete coefficient step into
the discrete countable module $Q$, including on compact parameter
families.  It is equivariant for the actual group actions on the
support and compact-dual models.  All these assertions allow
arbitrary root denominators in $W$.
\end{proposition}

\begin{proof}
First work at a bounded root level $h$.  In the normalized whole
basis, multiplication by \cref{eq:cutoff-trace-matrix} converts
\cref{eq:cutoff-residue-pairing} into the scalar residue test on each
basis coefficient.  Testing against all elements of
$x_h^{-D_0}A^{1/h}$ detects every strictly negative-y monomial of this
product whose x-order is at most $D_0/h$.  There is no undetected
smaller y-pole at this level: every positive $J_j$ in the root grid is
at least $1/h$.  Thus
\begin{equation}
 \lambda_h(w)=0
 \quad\Longrightarrow\quad
 \ord_x(w)>\frac{D_0-c}{h}
 \quad\text{in the normalized basis at level }h.
 \label{eq:cutoff-zero-test}
\end{equation}
Here one multiplies the detected product by $\mathcal T_h^{-1}$,
which loses at most $c/h$ in x-order.  These are invertible matrices
on the negative-y quotient: their entries and those of their inverses
have no negative y-powers.  Consequently projection to negative
y-powers does not invalidate their inverse identity.  This justifies
using the inverse matrix in \cref{eq:cutoff-zero-test}.

Conversely, suppose that $w$ has x-order at least $\epsilon>0$ in
the normalized basis at level $h$.  At a later level $h'$ the bound
is still at least $\epsilon$.  The trace/volume product then has
x-order at least $\epsilon-c/h'$.  If
$h'>(D_0+c)/\epsilon$, it pairs to zero with
$\mathcal P^{1/h'}$.  In particular every bounded-root element of
$W_+$ maps to zero in $Q$.  If a bounded-root element maps to zero in
$Q$, its functional is zero at some later level of the direct system;
\cref{eq:cutoff-zero-test,eq:cutoff-residue-margin} then put it in
$W_+$.  We have proved
\begin{equation}
 \ker(\pi_{\mathrm{fin}})=W_{\mathrm{fin}}\cap W_+.
 \label{eq:cutoff-finite-kernel}
\end{equation}

To prove surjectivity, take a functional at any stage of
\cref{eq:cutoff-Q}.  On the whole free lattice its finite coefficient
expression is represented by finite scalar principal parts.  Applying
the inverse trace/volume matrix gives a bounded-root representative
$v$ for \cref{eq:cutoff-residue-pairing}; it lies in the whole Laurent
module because that matrix has a fixed x-pole bound and no negative
y-powers.  Replace $v$ by $ev$.  These two functionals agree after one
transpose transition: the image of the corresponding transfer lies
in $eB'$, where multiplication by $e$ is the identity.  Thus $ev$
represents the same element of the colimit and belongs to
$W_{\mathrm{fin}}$.  This argument uses the whole trace matrix and
only then projects by $e$; it does not specialize an idempotent to
the boundary or presume a free model for its component there.

We next extend the map to $W$.  Express $w\in W$ on one fixed whole
normalized basis.  Truncate to the finitely many terms of x-order
$I\leq M$, obtaining $v_0$.  Let $c_e$ be a fixed x-pole bound for
the matrix of $e$ in that basis.  Taking $M>c_e+1$ gives
\[
 v=ev_0\in W_{\mathrm{fin}},\qquad
 w-v=e(w-v_0)\in W_+.
\]
The finitely many terms of $v_0$ have a common root denominator, and
$e$ has a bounded root denominator at this fixed basis level, so the
first inclusion holds.  Define $\pi(w)=\pi_{\mathrm{fin}}(v)$.
Any two such choices differ by a bounded-root element of $W_+$;
\cref{eq:cutoff-finite-kernel} proves independence of the choice.
If $\pi(w)=0$, the bounded-root representative $v$ belongs to $W_+$,
so $w$ also belongs to $W_+$.  This proves exactness.

For continuity, fix a complete weighted step and a normalized basis.
A sufficiently small norm ball has all its nonzero monomials of
weight greater than some $a>0$.  Such monomials necessarily have
$I>a/L$, because $w_{L,\delta}(I,J)\leq LI$ when $I\geq0$ and is
nonpositive when $I<0$.  The ball is therefore contained in $W_+$.
The kernel of the linear map to discrete $Q$ is open in this step.
The same argument in supremum norms works for a compact parameter
space: only finitely many bounded-x monomials remain uniformly in
the parameter, and each retained $\kk$-coefficient is locally
constant.  Its image in $Q$ is finite.

It remains to identify the group action, rather than just the kernel
as a vector space.  Work first at finite root level, and put
$A_h=\kk[[x_h,y_h]]$, $\mathfrak m_h=(x_h,y_h)$, and
$\Omega_h=\widehat\Omega^r_{A_h/\kk}$.
Write $\alpha=\Tr(wp_h\eta_h)$ as an ordinary top form, including
the factor $x_h^{-1}$ from $d\log x_h$, with its negative-y
\v Cech convention and its ordered-coordinate signs.
Choose $N$ so that $x_h^N\alpha$ has nonnegative x-orders.  Its
truncation modulo $x_h^N$ is a finite principal-part class
\[
 \xi_N\in H^{r-1}_{(y_h)}(\Omega_h/x_h^N\Omega_h)
       =H^{r-1}_{\mathfrak m_h}(\Omega_h/x_h^N\Omega_h).
\]
The equality holds because $x_h$ is nilpotent on the quotient.
The connecting map for multiplication by $x_h^N$ on $\Omega_h$
sends this class to the all-negative ordinary principal part of
$\alpha$ in $H^r_{\mathfrak m_h}(\Omega_h)$, with the coordinate
order fixing the \v Cech sign.  This identifies the finite
truncation with an intrinsic local-cohomology class.  If
$g(x_h)=u x_h$, the naturality square has $g$ on the source copy of
$\Omega_h$ and $u^{-N}g$ on the middle copy.  Consequently the
connecting map sends $u^{-N}g\xi_N$ to the $g$-translate of the
class of $\alpha$, as required.

These maps are the maps on the finite support \v Cech model of
\cref{prop:laurent-model}.  In particular a translation
$y_h\mapsto y_h+c(x_h)$ gives finite generalized-fraction
expansions modulo $x_h^N$, since $c(x_h)$ is nilpotent there.
The two-ended x-convention $(f_+,-f_-)$ ensures that multiplication
crossing x-order zero is ordinary Laurent multiplication.
For a convergent bounded-root expression, truncation above all fixed
x-pole bounds, including the group basis matrix bounds, changes
neither residue before or after the given action.  Thus the same
naturality applies to every expression being paired.

The generalized-fraction transformation law and the independence of
regular parameters for formal residue on
$H^r_{\mathfrak m_h}(\Omega_h)$ hold also in characteristic $p$
\citep[Lemmas~(7.2)(b) and~(7.3), pp.~60--67]{lipmanResidues1984}.
Here $A_h$ is a complete regular local algebra over the perfect
field $\kk$ with residue field $\kk$, so all hypotheses of that
result hold.  The finite étale trace is natural, the action preserves
$(x)$ and the whole lattice, and the transported volume is invariant
on the $e$-component.  These facts give
\[
 \lambda_h(gw)(p_h)=\lambda_h(w)(g^{-1}p_h)
 \qquad(g\in G).
\]
This proves equivariance on bounded-root expressions and is compatible
with \cref{eq:cutoff-residue-transition}.  Finite-root expressions
are dense after a complete-step enlargement.  Both the group action
and $\pi$ are continuous into enlarged steps and the discrete
quotient by \cref{lem:laurent-action} and the preceding paragraph.
Passing to the limit proves equivariance on all of $W$.  In particular
$W_+$ is $G$-stable.  Its positive complete steps also satisfy the
uniform action condition of \cref{def:cts-steps}.  Indeed, for a step
with x-gap $\epsilon$, pass to a later normalized basis $h'$.
The scalar substitutions preserve x-order by
\cref{eq:laurent-action-order}; the remaining group basis matrices
have one uniform pole bound $c_G/h'$, because they are root transports
of the matrices on the original whole model.  Choose $h'$ with
$c_G/h'<\epsilon/2$.  The entire compact group family then has x-gap
at least $\epsilon/2$, and \cref{lem:laurent-action} puts it
continuously into one common complete positive step.  The same
argument applies to compact parameter families.  Compact-envelope
integration from \cref{lem:cts-comparison} therefore defines the
completed group-ring action on $W_+$.  The map $\pi$ commutes with
that action: approximate a distribution by finite group-ring sums
in a compact envelope and use continuity of $\pi$ into $Q$.

This proof uses finite principal-part residue
invariance; it does not presume a separate duality theorem for
completed functions at infinite root level.
\end{proof}

\subsection{Contraction and an open boundary subgroup}

The residue sequence already makes the homology of its kernel
countable. Indeed, take a resolution by finite free completed group-ring
modules in every degree, as in \cref{lem:cts-comparison}. Its coefficient
chain complex has terms $V^{a_i}$, with each $a_i$ finite, for a
coefficient module $V$. Thus \cref{eq:cutoff-residue-sequence} gives a
termwise exact sequence of chain complexes. The complex for the
countable module $Q$ is degreewise countable, whereas $H_i(G,W)$ is
finite by \cref{cor:laurent-finite}. The long exact sequence gives
\begin{equation}
 H_i(G,W_+)\text{ is countable for every }i.
 \label{eq:cutoff-positive-countable}
\end{equation}

We must strengthen this to vanishing. On a complete positive coefficient
step, the cycles form a Polish group, and the subgroup consisting of
actual boundaries has countable index by
\cref{eq:cutoff-positive-countable}. We will show that this subgroup
is analytic, so a Baire-category argument makes it open. Frobenius
contraction then puts an iterate of every cycle in that open subgroup;
inverse Frobenius supplies a primitive for the original cycle. We first
establish the contraction and the required openness criterion.

\begin{lemma}[Contraction in one complete positive step]
\label[lemma]{lem:cutoff-contraction}
For every integer $\nu\geq0$ and every tuple $w\in W_+^\nu$, there
are an integer $a_*\geq0$, a complete weighted step written in the
original whole basis, and
$\epsilon,\alpha>0$ such that all $\Phi^{a_*+a}w$ belong to the
closed subspace with x-order at least $\epsilon$, and
\begin{equation}
 \|\Phi^{a_*+a}w\|\leq
       \exp(-\alpha h_0^a)\quad(a\geq0).
 \label{eq:cutoff-contraction}
\end{equation}
The statement applies without a bound on the root denominators of
the coefficient series.
\end{lemma}

\begin{proof}
The empty tuple is immediate. Otherwise, since the tuple is finite,
choose one normalized basis at level $h_0^{a_*}$ on which all
components have a common positive x-gap.  Applying $\Phi^{a_*}$ clears this
\emph{basis} denominator.  It need not clear the root denominators
in the series.  In the original basis the resulting tuple $v$ still
has a positive x-gap, since all multiplication laws in $A'$ are
integral.

Choose a common weighted step $(L_1,\delta)$ for the finitely many
series of $v$, with $I\geq\epsilon>0$. Their original weights are
bounded below.
Increasing $L_1$ to $L$ adds $(L-L_1)I$ to every weight.  Choose $L$
large enough that all nonzero terms now have weight at least
$\alpha>0$.  They still form null series in this new step.  Under
$h_0^a$-powering the coefficient exponents are multiplied by $h_0^a$,
so these weights are at least $\alpha h_0^a$.  Each
$b_i^{h_0^a}$ expands on the original basis with coefficients in
$A$.  These coefficients add nonnegative x- and y-powers;
projection to the strictly negative-y part can only discard terms
or improve the weight.  Their coefficients have norm at most one.
This proves \cref{eq:cutoff-contraction}.  The x-gap is multiplied
by $h_0^a$, hence remains at least $\epsilon$ in the same complete
step.
\end{proof}

\begin{lemma}[A countable-index analytic subgroup]
\label[lemma]{lem:cutoff-polish}
Let $Z$ be a Polish topological group and $B\subseteq Z$ an analytic
subgroup in the sense of descriptive set theory: $B$ is a continuous
image of a Polish space. If $Z/B$ is countable, then $B$ is open.
\end{lemma}

\begin{proof}
An analytic subset of a Polish space has the Baire property
\citep[Theorem~21.6]{kechris1995}.  If $B$ were meagre, its countably
many cosets would make $Z$ meagre in itself, contrary to the Baire
theorem.  Thus $B$ is nonmeagre.  Pettis's theorem
\citep[Theorem~9.9]{kechris1995} says that the product of a
nonmeagre set with the Baire property and its inverse contains a
neighborhood of the identity.  For a subgroup this product is $B$
itself, proving the assertion.
\end{proof}

\begin{lemma}
\label{lem:cutoff-positive-acyclic}
For the completed-resolution group homology of
\cref{sec:continuous},
\[
 H_i(G,W_+)=0\qquad(i\geq0).
\]
\end{lemma}

\begin{proof}
Use the finite free resolution and the countability statement
\cref{eq:cutoff-positive-countable} above. We first construct the
complete cycle spaces to which \cref{lem:cutoff-polish} applies.
Take a normalized basis level, rational
weights $L>\delta>0$, and a rational positive cutoff
$\epsilon>0$.  In the whole weighted null-series space require
x-order at least $\epsilon$ and the equation $ew=w$.  The first
condition is a closed coordinate condition; the second is closed
because $e$ and the identity are continuous into a common enlarged
complete step.  This gives a complete separable metrizable
$\kk$-vector space $E\subset W_+$.  Separability follows from the
countable exponent index and finite coefficient field.  These
spaces form a countable collection covering $W_+$: the basis levels
are countable, and rational weights and positive cutoffs can be
chosen cofinally.  Each inclusion between two specified coefficient
steps is continuous into some common enlarged step.

Fix $i$ and one such $E$.  The cycles
\[
 Z_E=\{z\in E^{a_i}: dz=0\}
\]
form a closed subspace of a Polish group.  In fact the differential
is continuous from this finite product into a common complete step
by \cref{lem:laurent-action,lem:cts-comparison}.  Let
\[
 B_E=Z_E\cap d(W_+^{a_{i+1}}).
\]
Its index in $Z_E$ is countable by
\cref{eq:cutoff-positive-countable}.  It is analytic for the
following explicit reason.  For each complete positive primitive
step $E'$, the set
\[
 \{(z,u)\in Z_E\times(E')^{a_{i+1}}:du=z\}
\]
is closed: compare the two sides in a common enlarged complete
step.  Projecting this closed subset of a Polish product onto
$Z_E$ gives an analytic set.  There are countably many choices of
$E'$, and their projected union is exactly $B_E$.  Thus $B_E$ is
an analytic subgroup, and \cref{lem:cutoff-polish} makes it open.
This argument does not require the global boundary image in $W_+$
to be closed or the inductive topology to be strict.

Now let $z$ be any cycle in $W_+^{a_i}$.  Absolute Frobenius is
$\kk$-linear and $G$-equivariant.  It commutes with the completed
group-ring differential: this holds for finite group-ring sums and
then for their convergent actions on the complete steps.  Apply
\cref{lem:cutoff-contraction} and choose the rational weights and
cutoff within its bounds.  The forward iterates after the initial
basis clearing lie in one $Z_E$ and tend to zero there.  Since
$B_E$ is open, some iterate $\Phi^a z$ is $du$ with
$u\in W_+^{a_{i+1}}$.  Both $\Phi$ and its inverse preserve $W_+$
and commute with $d$.  Hence $z=d(\Phi^{-a}u)$.  Every cycle
bounds, as required.
\end{proof}

\begin{theorem}[Finite stable transfer image]
\label{thm:cutoff-stable}
For every $i\geq0$ the compact vector space
$M_i=H^i(G,\mathcal P)$ satisfies
\begin{equation}
 I_i:=\bigcap_{a\geq0}S^aM_i\quad\text{is finite}.
 \label{eq:cutoff-stable}
\end{equation}
No uniform bound on its size as the prime, height, degree, or finite
marking level varies is asserted.
\end{theorem}

\begin{proof}
The preceding acyclicity and
\cref{eq:cutoff-residue-sequence} identify $H_i(G,Q)$ with
$H_i(G,W)$, hence make it finite.  Compact/discrete duality and the
finite free resolution give
\[
 H_i(G,Q)
   =\colim_{S^\vee}H_i(G,\mathcal P^\vee)
   =\colim_{S^\vee}M_i^\vee.
\]
Filtered colimits commute with kernels and cokernels of vector-space
maps, so no derived interchange is involved here.  Dualizing gives
\[
 H_i(G,Q)^\vee=\varprojlim_S M_i.
\]
This inverse limit is finite.  Its projection to the first term has
image $I_i$.  To check surjectivity onto that image description, note
that every $S^aM_i$ is compact.  For $x\in I_i$, the nonempty compact
sets
\[
 \{y\in S^aM_i:Sy=x\}
\]
are nested.  Their intersection gives a preimage in $I_i$.
Thus $S:I_i\to I_i$ is surjective, and successive choices of
preimages in $I_i$ lift $x$ to the inverse limit.  It follows that
$I_i$ is finite.  We have used a surjection from the inverse limit
onto the stable image; we have not assumed them equal before proving
finiteness.
\end{proof}

\subsection{Compact power-series modules and localization}

The stable image is now finite in compact lattice cohomology. To pass
to the intersection of its images in localized cohomology, we will need
two additional inputs: next-height finiteness, which controls boundary
strips, and countability of a consecutive transfer-image quotient.
The next lemma isolates the compact algebra used in this passage.

\begin{lemma}
\label[lemma]{lem:cutoff-compact-operator}
Let $M$ be a profinite $\kk$-vector space and let $S:M\to M$ be a
continuous linear map.  Assume that
$I=\bigcap_{a\geq0}S^aM$ and $M/SM$ are finite.  Then $S$ is an
automorphism of $I$, and $V=M/I$ is a finitely generated
$\kk[[z]]$-module with $z$ acting by $S$.  Its original topology is
the z-adic topology.  The locally nilpotent subgroup
\[
 T_S(M)=\bigcup_{a\geq0}\ker(S^a:M\to M)
\]
is finite and closed.  In particular $\ker S$ is finite.
\end{lemma}

\begin{proof}
The compact preimage argument in the preceding proof gives
$S(I)=I$, and finiteness makes this restriction invertible.  In
$V=M/I$ the nested compact subspaces $S^aV$ have zero intersection.
They shrink \emph{uniformly} to zero: for every neighborhood $O$
of zero there is an $a_0$ such that $S^aV\subset O$ for every
$a\geq a_0$.  Otherwise the compact sets
$S^aV\cap(V\setminus O)$ would have nonempty intersection.

For $f=\sum_{j\geq0}c_jz^j$ define
\[
 f v=\lim_{t\to\infty}\sum_{j=0}^{t-1}c_jS^jv.
\]
Every sufficiently late finite tail belongs to an arbitrarily
specified open linear subspace, uniformly for $v\in V$.
Completeness therefore supplies the limit.  Polynomial module
identities pass to these uniform limits and give a continuous
$\kk[[z]]$-action.

Choose lifts $v_1,\ldots,v_s$ of a basis of the finite space
$V/SV=M/SM$.  Given $v=v^{(0)}$, recursively choose coefficients
$c_{ij}\in\kk$ and $v^{(j+1)}\in V$ such that
\[
 v^{(j)}=\sum_{i=1}^s c_{ij}v_i+Sv^{(j+1)}.
\]
After $t$ steps the error is $S^tv^{(t)}$, which tends uniformly
to zero.  Thus
$v=\sum_i(\sum_jc_{ij}z^j)v_i$; the continuous map
$\kk[[z]]^s\to V$ is surjective.  Each successive quotient
$S^jV/S^{j+1}V$ is an image of $V/SV$.  Hence every $V/S^aV$ is
finite, and its compact kernel $S^aV$ is open.  Uniform shrinking
makes these kernels a neighborhood basis, proving the topology
assertion.

The Noetherian module $V$ has a stabilizing ascending sequence
$\ker z\subseteq\ker z^2\subseteq\cdots$.  Its union is a
finitely generated submodule killed by some $z^t$, hence finite over
the finite ring $\kk[[z]]/(z^t)$.  The map from $T_S(M)$ to this
union is injective, since a locally nilpotent element of $I$ is
zero.  Therefore $T_S(M)$ is finite and closed.  In fact this map
is surjective: if $S^a x\in I$, subtract the unique element of $I$
with the same $S^a$-image to obtain a nilpotent lift.  The contraction
and the power-series action in this argument concern $M/I$; they
are not asserted on the finite summand on which $S$ is invertible.
\end{proof}

Return to $M_i=H^i(G,\mathcal P)$ and set
\[
 Y_i=H^i(G,B'),\qquad X_i=H^i(G,B_U).
\]
Let $j_i:M_i\to Y_i$ be localization and let
$f_i:M_i\to X_i$ be localization followed by projection by $e$.
For the remainder of this section assume the next-height finiteness
hypothesis of \cref{lem:mark-strips}.  That lemma says that $j_i$
has countable kernel and cokernel for every fixed lattice shift.
It need not say this about $f_i$ before taking a positive transfer
image: the complementary component of $B'$ must first be removed.

For every $a\geq1$ there is a map
\begin{equation}
 f_{i,a}:S^aM_i\longrightarrow S^aX_i
 \quad\text{with countable kernel and cokernel}.
 \label{eq:cutoff-image-comparison}
\end{equation}
Here is the verification.  The map
$S^aM_i\to S^aY_i$ has kernel contained in $\ker j_i$.
Its cokernel is an image of $Y_i/j_i(M_i)$ under $S^a$, and is
therefore countable.  For $a\geq1$, $S^aY_i$ lies in the
$e$-summand, where projection identifies it with $S^aX_i$.
This proves \cref{eq:cutoff-image-comparison} without imposing a
countable-kernel assertion on the uniterated projection.

\begin{proposition}[The compact localized intersection]
\label[proposition]{prop:cutoff-compact}
Fix $i\geq0$ and $N\geq1$.  Suppose
\begin{equation}
 S^NX_i/S^{N+1}X_i\quad\text{is countable}.
 \label{eq:cutoff-consecutive-hypothesis}
\end{equation}
Then $S^NM_i/S^{N+1}M_i$ is finite, and
$(S^NM_i)/I_i$ is a finitely generated $\kk[[z]]$-module with
$z=S$.  The operator $S$ on $S^NX_i$ has countable kernel and
cokernel.  Moreover $\ker(f_i|_{S^NM_i})$ is finite and closed,
$f_i$ is injective on $I_i$, and
\begin{equation}
 \bigcap_{a\geq0}\im(S^aM_i\longrightarrow X_i)
      =f_i(I_i)\cong I_i.
 \label{eq:cutoff-localized-intersection}
\end{equation}
In particular the intersection is finite.  No topology on $X_i$
and no closed-image assertion for its cochain differential is
required.
\end{proposition}

\begin{proof}
Put $A_N=S^NM_i$, $Z_N=S^NX_i$, and $f=f_i|_{A_N}$.
The induced map
\[
 A_N/SA_N\longrightarrow Z_N/SZ_N
\]
has countable kernel.  Indeed, the subgroup represented by
$\{a\in A_N:f(a)\in SZ_N\}$ maps to
$SZ_N/f(SA_N)$, which is countable by
\cref{eq:cutoff-image-comparison} at $N+1$; the kernel of this
map is an image of $\ker f$, countable by the comparison at $N$.
Its image in $Z_N/SZ_N$ is countable by
\cref{eq:cutoff-consecutive-hypothesis}.  Thus $A_N/SA_N$ is
countable.  It is a compact Hausdorff group, because $SA_N$ is a
compact image and hence closed.  A countable compact Hausdorff
group is finite: its countable cover by closed singletons, together
with Baire, forces a singleton to be open; the group is then
discrete and compact.  This proves finiteness of $A_N/SA_N$.

The stable intersection for $S$ acting on $A_N$ is $I_i$.
Apply \cref{thm:cutoff-stable,lem:cutoff-compact-operator}.
This gives the asserted finite generation and makes both
$\ker(S|_{A_N})$ and $T_S(A_N)$ finite.

We first obtain countability of the kernel on $Z_N$ without using
closedness of $\ker f$.  Write $Z_0=f(A_N)$ and
$C_N=Z_N/Z_0$, which is countable.  An element of
$\ker(S|_{Z_0})$ lifts to an $a\in A_N$ with
$Sa\in\ker f$.  The latter kernel is countable and each
nonempty fiber of $S:A_N\to A_N$ is finite.  Consequently
$S^{-1}(\ker f)$, and hence $\ker(S|_{Z_0})$, is countable.
Mapping $\ker(S|_{Z_N})$ to $C_N$ proves its countability.
The cokernel is countable already by
\cref{eq:cutoff-consecutive-hypothesis}.

It remains to prove the stronger assertion about $\ker f$, for
which countability alone would not suffice.  Let
$\alpha\in\ker f\subseteq S^NM_i$.  Choose a
$\mathcal P$-valued cocycle $c$ such that $S^Nc$ represents
$\alpha$.  Since $N\geq1$, this representative is already
$e$-supported after localization.  If $i>0$, vanishing in $X_i$
supplies a bounded-pole primitive $b$ with
\[
 db=S^Nc
\]
in the cochain complex for $B_U$, viewed inside that for $B'$.
The convention of \cref{lem:cts-exact,lem:cts-comparison} means
that this one primitive belongs to some whole lattice
$x^{-D}A'$ uniformly on its compact cochain parameters.
By \cref{eq:cutoff-eventual-integrality}, all sufficiently high
$S^ab$ are $\mathcal P$-valued.  The chain-map identity gives
$d(S^ab)=S^a(S^Nc)$ there.  Thus $S^a\alpha=0$ in $M_i$.
When $i=0$, vanishing of the already $e$-supported localized cocycle
means that the cocycle itself is zero, since
$\mathcal P\hookrightarrow B'$ is injective.  This proves local
$S$-nilpotence of $\ker f$ in every degree.  The exponent may
depend on its chosen bounded-pole primitive; no uniform exponent
was assumed.

We now have $\ker f\subseteq T_S(A_N)$, so $\ker f$ is finite
and closed.  Its intersection with $I_i$ is zero, because $S$
acts invertibly on $I_i$.  For an element $\xi$ belonging to all
the images in \cref{eq:cutoff-localized-intersection}, discard
the first $N$ terms and consider
\[
 E_a=\{v\in S^aA_N:f(v)=\xi\},\qquad a\geq0.
\]
These sets are nonempty and nested.  Each is a coset of
$\ker f\cap S^aA_N$, hence finite and closed in the compact
space $A_N$.  Their intersection is nonempty.  Any point in it
belongs to $I_i$ and maps to $\xi$.  This proves the nontrivial
inclusion in \cref{eq:cutoff-localized-intersection}; the reverse
inclusion is immediate.  It also explains why no Hausdorff topology
on localized cohomology was used.
\end{proof}

The sequence of implications is now available at every fixed allowed
marking level.  The finite stable image comes from the analytic
\emph{transpose} quotient, whereas the final compact argument
concerns the original bounded-pole cochains.  In particular none of
the arguments in this section identifies incomplete ordinary
cohomology on the open with the ordinary cohomology of its analytic
completion.
